\begin{textsample}{POS Dim 1 – mt – Score 43.00 – mt/mt\_14.txt}  \label{ex:f1_pos_080}
4.
AS DIVERSIDADES EDUCACIONAIS NA EDUCAÇÃO BÁSICA Neste \textbf{tópico}, apresentamos diretrizes relativas às diversidades educacionais na educação básica, a serem contempladas tanto no currículo de escolas específicas quanto no currículo de escolas urbanas de ensino regular. 4.1.
A Educação Especial na Perspectiva da Educação Inclusiva na Educação Básica Ao serem consideradas as questões da Educação Especial na Perspectiva da Educação Inclusiva, entende -se como \textbf{oportuna} a construção de \textbf{um} espaço reflexivo, de modo a dar visibilidade ao trabalho com as diferenças, valorizando a diversidade, sustentada em princípios éticos, políticos e estéticos, traçados pelas Diretrizes Curriculares Nacionais da Educação Básica, bem como, pela BNCC, avançando de forma sistemática na \textbf{direção} da construção de \textbf{uma} sociedade \textbf{justa}, democrática e inclusiva.
Desse modo, a educação inclusiva deixa de ser \textbf{uma} possibilidade e passa fazer parte da rotina educacional do nosso estado, \textbf{uma} vez \textbf{que} \textbf{implica} repensar a escola em seu papel educacional e social, no sentido de assegurar a todos os estudantes, o direito ao desenvolvimento de competências \textbf{que} lhe assegurem o direito de aprender.
Em conformidade com os textos legais, o princípio básico \textbf{que} orienta a Educação Especial na perspectiva inclusiva é o de \textbf{que} todas as crianças devem aprender juntas, independentemente de quaisquer dificuldades ou diferenças físicas, intelectuais, sociais, emocionais, linguísticas ou outras.
Esse conceito é a base de sustentação de \textbf{uma} só escola para todos.
Além de trabalhar o conhecimento de forma \textbf{sistematizada}, a escola deve objetivar processos de aprendizagem de acordo com as necessidades básicas de cada estudante.
Nesse sentido, a Educação Especial, conforme definida pelos dispositivos legais pertinentes, é \textbf{uma} modalidade educativa configuradora de proposta pedagógica \textbf{que} hospeda concepções, formas de exequibilidade e serviços especializados, institucional e operacionalmente estruturados, assentados para apoiar, complementar e suplementar o processo de educação escolar e o itinerário formativo, em nível ascendente, de estudantes com deficiência, transtorno global do desenvolvimento e altas habilidades ou superdotação.
Esse conceito envolve o princípio democrático de educação para todos, inserida na transversalidade das diferentes etapas e modalidades da educação escolar.
Desse modo, estamos diante de \textbf{uma} noção \textbf{que} ultrapassa o conceito de escola especial.
Para essas \textbf{tarefas}, os professores devem buscar formação, de maneira permanente, objetivando atuar na perspectiva de \textbf{uma} sala de aula \textbf{que} já não \textbf{foca} a deficiência do estudante, mas o tipo de mediação pedagógica, resposta educativa e de recursos e apoios \textbf{que} a escola disponibiliza para \textbf{que} este estudante obtenha sucesso escolar.
A resposta para a internalização curricular pelo \textbf{educando} com deficiência, transtornos globais do desenvolvimento e altas habilidades ou superdotação está nas formas e nas condições de aprendizagem \textbf{que} lhe são oferecidas.
De partida, \textbf{convém} \textbf{atentar} para o fato de \textbf{que} estudantes com deficiência visual, auditiva ou física detêm, em princípio, todas as condições para o acompanhamento dos colegas de turma.
Esse não é o caso daqueles com deficiência intelectual, paralisia cerebral, autismo e situações cognitivas diferenciadas, \textbf{que} \textbf{necessitam}, de fato, de \textbf{um} plano de ensino com configuração específica, com objetivos e organização didáticos diferenciados.
Aspectos \textbf{que}, igualmente, devem ser considerados no caso de estudantes com altas habilidades.
Na \textbf{direção} \textbf{ora} indicada, \textbf{uma} possibilidade se impõe \textbf{que} é o avanço no fortalecimento de \textbf{um} processo diferenciado de \textbf{ensino-aprendizagem} com foco na progressão escolar dos alunos.
De fato, o princípio da aprendizagem diferenciada, decisivo nas práticas escolares inclusivas, não deve caminhar para \textbf{uma} categorização dos estudantes.
A questão \textbf{aqui} é de natureza totalmente diversa.
Tem a \textbf{ver}, sob o ponto de vista pedagógico, com currículo e, sob o ponto de vista da gestão escolar, com procedimentos.
Mas o currículo definido não impede sua execução com suficiente flexibilidade para \textbf{que} se adequem aos estudantes na conformação de suas necessidades.
Quanto aos procedimentos, a escola, em sua estrutura organizacional, deve permitir \textbf{que} os estudantes com deficiência, transtornos globais do desenvolvimento e altas habilidade ou superdotação, possam participar maximamente dos ambientes escolares regulares.
Como aponta Coll ( 1995, p. 301 ) “ (... ) é importante \textbf{que} os estudantes com deficiência, transtornos globais do desenvolvimento e altas habilidade ou superdotação participem de \textbf{uma} programação tão normal quanto possível e tão específica quanto suas necessidades requeiram ”.
E acrescenta \textbf{que} essa disposição institucional \textbf{implica} em “ (... ) dispor de procedimentos e modelos de adequação individualizada do currículo \textbf{que} sirvam para assegurar este difícil e \textbf{imprescindível} equilíbrio ”.
No currículo da educação inclusiva e de \textbf{uma} só escola para todos, o professor da Educação Especial são todos os professores no sentido de \textbf{que} os docentes das classes regulares devem estar preparados para atender e acompanhar os estudantes com deficiência, transtornos globais do desenvolvimento e altas habilidade ou superdotação e ajudá -los a construir \textbf{uma} autoimagem positiva, \textbf{uma} visão de mundo real e se aceitar como diferente.
Nesse rumo, a proposta pedagógica da escola, base de referência dos contornos da prática escolar, deve identificar, no currículo em ação, \textbf{um} instrumento ( recurso operativo ) para o desenvolvimento dos estudantes.
Para tanto, a escola trabalha com \textbf{um} planejamento situacional, envolvendo aspectos relevantes, dentre os quais cabe destacar : atitude permeável à diversificação e à flexibilização dos processos de ensino ; adoção de currículos abertos e de ofertas curriculares múltiplas em sua execução, em substituição a estruturas curriculares uniformes, homogeneizantes e inflexíveis ; inclusão, em todas as escolas e sempre \textbf{que} necessário, de professores especializados, de variados serviços de apoio e de outros não-convencionais, indispensáveis à viabilização do processo educacional inclusivo.
Em síntese, como destaca Peter Mittler ( 2003, p. 183 ), é necessário “ preparar todos os professores para ensinar a todos os alunos ”.
Mais do \textbf{que} isso, desenvolver todos da escola para \textbf{que} possam ajudar no desenvolvimento de todos \textbf{que} estão na sala de aula.
Esses objetivos estão calçados pelos princípios fundamentais da \textbf{dignidade} da pessoa humana, da busca da identidade e do exercício da cidadania.
Obviamente, a Educação Especial busca realizar seus objetivos por meio de campos \textbf{metodológicos} e práticas pedagógicas especiais, \textbf{explicitadas} em alternativas de atendimentos diferenciados e recursos humanos devidamente especializados.
Esses meios de operacionalização do acesso ao currículo escolar devem servir de parâmetro de referência para o planejamento das atividades docentes e representam ajustes tênues da sala de aula.
Rebatem, também, no quadro de competências e habilidades a serem desenvolvidas, bem como nos aspectos de tempos e espaços de aprendizagem e nos procedimentos de avaliação.
Há \textbf{uma} certa convergência de entendimento quanto ao estudante da Educação Especial.
A literatura especializada o define de forma mais ou menos consensual, embora seja possível vislumbrar, \textbf{aqui} e ali, \textbf{percepções} com tendências a ressaltar alguns traços prevalecentes.
Dentro dessa ampla e convergente compreensão, esse estudante tem sido definido como detentor de mais de \textbf{uma} das seguintes características : - Dificuldades acentuadas e duradouras de aprendizagem ; - Limitações significativas no processo de seu desenvolvimento intelectual, com reflexos diretos ou indiretos no comportamento diário e no sequenciamento permanente das atividades curriculares programadas intensiva e extensivamente ; - Dificuldades ou limitações podem ser de gênese orgânica ou, ainda, podem estar vinculadas a condições existenciais sob a forma de disfunções, restrições ou deficiências mais \textbf{expressivas} ; - Limitações comunicacionais ; - Necessidade de uso de formas de sinalização diversas dos demais estudantes, com exigência de apoio em linguagens e códigos correspondentes ; - Alta capacidade de domínio de componentes de aprendizagem, expressa no encurtamento do tempo para a apropriação de conceitos, procedimentos, e para o desenvolvimento de habilidades e atitudes resolutivas, bem como na captura performática de espaços múltiplos de campos de conhecimento.
Em qualquer tempo, o estudante da Educação Especial, ingressante no processo de escolarização contará com o apoio do Serviço Educacional Especializado, de modo a dar suporte na avaliação de suas condições cognitivas e de domínio dos conhecimentos curriculares como requisito para a definição de sua situacionalidade no nível de ensino e a sua fase no ciclo.
No Sistema de Ensino de Mato Grosso, os alunos com deficiência, transtornos globais de desenvolvimento e altas habilidades ou superdotação contam, adicionalmente, com formas de complementação ou suplementação ao processo de escolarização com duas plataformas de atendimento.
São elas : - Os Serviços de Apoio Pedagógico Especializado : são desenvolvidos em Classes Comuns ; Salas de Recursos Multifuncionais ; - Os Serviços Especializados : são desenvolvidos em Escolas Especializadas e em Escolas Especializadas Comunitárias, Confessionais ou Filantrópicas sem fins lucrativos.
Esses serviços possuem amplitude crescente, ficando sua configuração na dependência da disponibilidade de recursos institucionais, humanos, materiais e simbólicos em nível dos respectivos contextos ; - Classe Comum : espaço da escola regular, organizado para acolher todos os estudantes e, dessa forma, ensejar a convivência entre aqueles com deficiência, transtornos globais de desenvolvimento e altas habilidades ou superdotação e os demais ; - Salas de Recursos Multifuncionais : espaço flexível de aprendizagem, constituído por equipe multiprofissional e estruturado com duplo objetivo : realizar complementações, quando se tratar de atividades extensivas do currículo, ou \textbf{operar} suplementações, quando se tratar de atividades para o preenchimento de lacunas na aprendizagem.
Embora as instituições educacionais trabalhem com diferentes \textbf{abordagens} epistemológicas, a LDB, art. 3º, inciso III, denomina como pluralismo de ideias e de concepções pedagógicas, no campo da Educação Especial e recomenda \textbf{um} olhar cuidadoso nos pontos de vista gnosiológicos de Piaget e Vygotsky.
As contribuições desses pensadores são marcantes tanto no aspecto dos alinhamentos \textbf{teóricos}, quanto no \textbf{que} tange à fundamentação das práxis pedagógica.
A interlocução \textbf{dialógica} das ideias desses \textbf{autores} ajudará escolas e professores a construírem rotas de inclusão social e escolar com maior visibilidade.
O destaque das conexões profundas entre a dimensão cognitiva e a dimensão afetiva é a extraordinária contribuição do pensamento de Vygotsky.
Sobretudo, na formulação de intervenções psicopedagógicas no interior de \textbf{um} projeto de educação inclusiva e, portanto, de inserção de estudantes da Educação Especial, assim considerados \textbf{historicamente}, nas escolas regulares e nas salas de aula comuns.
Com as contribuições de Piaget e Vygotsky, é possível construir \textbf{uma} linha de fundamentação \textbf{teórica} no campo das práticas pedagógicas inclusivas.
Sobretudo, é possível potencializar as mediações adequadas, para a viabilização de \textbf{uma} Educação Especial não divorciada da ideia de \textbf{UMA} SÓ ESCOLA PARA TODOS.
Colocadas essas premissas, cabe ressituar a posição da Educação Especial, agora, na moldura da educação inclusiva, e ressaltar a transversalidade da modalidade.
Trata -se, portanto, de \textbf{uma} construção de escolaridade \textbf{que} caminha em sentido vertical quando potencializada nas finalidades e no corpo de objetivos e no currículo de cada nível de ensino e, em sentido horizontal, quando reforçada no corpo de metodologias e práticas pedagógicas no âmbito da própria modalidade educativa.
Essa constatação sinaliza a significativa evolução legal, \textbf{conceitual} e pedagógica por \textbf{que} passa a Educação Especial.
De fato, deixa de ser enxergada como \textbf{um} subsistema congelado na organização de \textbf{um} tipo de educação e passa a assumir \textbf{uma} compreensão sócio-educativa dinâmica, apoiada na concepção inclusivista de \textbf{uma} só escola para todos.
A escola de Educação Básica é o ambiente institucional próprio \textbf{que} a sociedade criou.
No processo de evolução humana, o indivíduo recebe informações \textbf{sistematizadas}, para reelaborá -las sob a forma de conhecimento e convertê -las em habilidades para ampliação de comunicação, produção de objetos, bens e serviços.
Tudo isso leva à formação de atitudes, valores e à construção cultural.
A Educação Básica é, portanto, o \textbf{alicerce} para a construção do saber organizado, no quadro de \textbf{uma} organização do conhecimento.
O texto \textbf{que} você tem em mãos foi desenvolvido a partir do entendimento de \textbf{que} as Competências Gerais apresentadas na BNCC não devem ser interpretadas como \textbf{um} componente curricular, mas segundo nos aclara o texto, de forma transdisciplinar, presentes em todas as áreas de conhecimento e etapas da educação.
Para tal, consideramos as competências, no texto da BNCC, como sendo a mobilização de conhecimentos, habilidades, atitudes e valores para resolver demandas da vida cotidiana, do exercício da cidadania e do mundo do trabalho.
Isso significa \textbf{que} competência é aquilo \textbf{que} permite aos estudantes desenvolverem plenamente cada \textbf{uma} das habilidades e aprendizagens essenciais apresentadas na Base.
No entanto, na perspectiva da educação inclusiva, é necessário \textbf{que} as instituições de ensino, bem como suas mantenedoras ofereçam apoio para \textbf{que} cada estudante desenvolva as competências necessárias para a realização dos seus projetos de vida, respeitado seu itinerário formativo.
Não se pode esquecer \textbf{que} a escola tem \textbf{uma} enorme responsabilidade no campo da socialização cívica e política.
Ela contribui para \textbf{qualificar} as interações e os valores socioculturais nessa fase decisiva de vida dos estudantes.
A Educação Básica trabalha o currículo com foco nas competências gerais trazidas pela BNCC e nas competências específicas dos diversos componentes curriculares e em seus campos de experiências, habilidades e objetos de conhecimentos.
Sendo assim, imperioso, portanto, ressituar o estudante da Educação Especial na escola regular, nas classes comuns e vinculadas aos níveis formais de ensino, imprimindo, a seu modo e dentro do seu padrão de subjetividade, as peculiaridades legais e pedagógicas de cada \textbf{um} deles.
É importante reconhecer o planejamento pedagógico e as intervenções docentes como os dois rumos prevalecentes da adequação curricular.
Sendo \textbf{que}, no âmbito da adequação curricular, a escola deverá definir clara e coletivamente : - O \textbf{que} o estudante deve aprender ( Aprendizagem significativa ) ; - Como o estudante deve aprender ( Procedimentos ) ; - Quando o estudante está apto a aprender ( Temporalidades ) ; - Tipologias seletivas para assegurar o processo eficiente de aprendizagem ( Formas de organização do ensino ) ; - Modalidades de avaliação e processos sequenciais de acompanhamento do progresso do estudante.
Essas decisões devem envolver o professor, o estudante e o coletivo de professores, \textbf{uma} vez \textbf{que} se trata de procedimentos de adequação encorpados na proposta pedagógica e no currículo em execução, no espaço da sala de aula.
As adequações curriculares são procedimentos processuais em \textbf{que} estudante e professor se envolvem com o objetivo de possibilitar ao aluno a superação de dificuldades de aprendizagem.
Desse modo, a adequação curricular é \textbf{uma} resposta da escola às necessidades dos estudantes com limitações de aprendizagem.
Nesse caso, o currículo se torna mais apropriado, sob o ponto de vista de \textbf{assimilação} dos estudantes.
O ponto de referência para as adequações é o currículo comum, geral, regular.
A partir dele, são adotadas formas mais confortáveis intelectualmente para o estudante progredir em sua escolaridade.
A adequação é resultado do conjunto de elementos materiais, de recursos pedagógicos e de formas de intervenção professor/ estudante, mobilizados para ampliar as chances de os estudantes com deficiência, transtornos globais do desenvolvimento e altas habilidades ou superdotação se desenvolverem em seu aprendizado.
Os procedimentos de avaliação são trabalhados com conceitos \textbf{baseados} em indicadores.
Estes, por sua vez, devem guardar relação intrínseca com as áreas de conhecimento.
A avaliação não é ocasional, mas diagnóstica e formativa, ou seja, desfoca -se dos \textbf{métodos} normativos e quantitativos e se hospeda numa moldura mais ampla e em processo mais dinâmico-coletivo onde se consideram os avanços do desenvolvimento pessoal, as articulações com o contexto e os ganhos qualitativos do \textbf{sujeito}.
Na visão da educação inclusiva, a avaliação reconceituada convoca a \textbf{singularidade} do \textbf{sujeito} à \textbf{centralidade} do processo educativo e da relevância pedagógica permanente, o \textbf{que} importa valorizar a \textbf{cognição} afetiva como fonte de interações das cadeias vitais e das redes de interdependência no âmbito de \textbf{uma} aprendizagem significativa.
Quando se aborda a questão da avaliação na Educação Especial, \textbf{posta} esta em posição de ressignificação porque pautada pelos parâmetros da educação inclusiva e, portanto, pelo olhar de modalidades emancipadoras de avaliação, é necessário compreender \textbf{que} a imposição epistemológica e cultural \textbf{que} exclui, da sala de aula comum, alguns tipos de estudantes ou os silencia, deve ser substituída por \textbf{uma} outra concepção pedagógica, por \textbf{uma} outra escola e por \textbf{um} outro professor.
A perspectiva deve ser reelaborada, em novas formas de a escola trabalhar, reconstruir, ressignificar e avaliar conhecimentos curriculares em ação.
Tem como ponto de partida o reconhecimento das imensas potencialidades da \textbf{singularidade} humana, com repercussão direta nas formas de ensinar, aprender e avaliar, \textbf{uma} vez \textbf{que} avaliar é colocar o \textbf{sujeito} em condições plenas de se \textbf{revelar} como identidade em construção.

% matched lemmas: abordagem, alicerce, aqui, assimilação, atentar, autor, basear, centralidade, cognição, conceitual, convir, dialógico, dignidade, direção, educar, ensino-aprendizagem, explicitar, expressivo, focar, historicamente, implicar, imprescindível, justo, metodológico, método, necessitar, operar, oportuno, ora, percepção, postar, qualificar, que, revelar, singularidade, sistematizar, sujeito, tarefa, teórico, tópico, um, uma, ver
\end{textsample}
